\documentclass[12pt,a4paper]{article}
\usepackage[UTF8]{ctex}
\usepackage{amsmath}
\usepackage{amssymb}
\usepackage{graphicx}
\graphicspath{{D:/code/LATEX/pscs_expmt/figures}}
\usepackage{geometry}
\usepackage{setspace}
\usepackage{indentfirst}
\usepackage{bm}
\usepackage{array}
\usepackage{url}
\usepackage{cite}
\usepackage{multicol}
\usepackage{fancyhdr}
\usepackage{booktabs}
\usepackage{titlesec}
\usepackage{caption}
\usepackage{lastpage}

% 页面设置
\geometry{top=2.5cm, bottom=2.5cm, left=2cm, right=2cm}
\raggedbottom

% 设置行距
\renewcommand{\baselinestretch}{1.25}

% 设置页眉页脚
\pagestyle{fancy}
\fancyhf{} % 清除所有页眉页脚
% 第一页特殊设置 - 只有顶部水平线，没有页眉线
\fancypagestyle{firstpage}{
    \fancyhf{} % 清除所有页眉页脚
    \renewcommand{\headrulewidth}{0pt} % 第一页没有页眉线
    \fancyfoot[C]{\zihao{-6}\songti 课程名称：\quad 实验日期：\quad 作者简介：} % 页脚
}
% 其他页设置 - 有页眉线
\fancyhead{} % 清除页眉
\renewcommand{\headrulewidth}{1pt} % 其他页有页眉线
\fancyfoot[C]{\zihao{-6}\songti 课程名称：\quad 实验日期：\quad 作者简介：} % 页脚

% 设置章节格式
\titleformat{\section}{\zihao{-4}\heiti}{\thesection}{1em}{}
\titleformat{\subsection}{\zihao{5}\heiti}{\thesubsection}{1em}{}
\titleformat{\subsubsection}{\zihao{5}\kaishu}{\thesubsubsection}{1em}{}

% 设置段前段后间距
\titlespacing*{\section}{0pt}{0.5\baselineskip}{0.5\baselineskip}
\titlespacing*{\subsection}{0pt}{0.5\baselineskip}{0.25\baselineskip}
\titlespacing*{\subsubsection}{0pt}{0.25\baselineskip}{0.25\baselineskip}

% 设置图表标题格式
\captionsetup[table]{font=small, skip=5pt}
\captionsetup[figure]{font=small, skip=5pt}

\begin{document}

% 第一页使用特殊页眉样式
\thispagestyle{firstpage}

% 顶部水平线（只保留这一条）
\vspace*{-2cm}
\hrule width \textwidth height 1pt
\vspace{1.5cm}

% 中文标题部分 - 通栏排版
\begin{center}
    {\zihao{-2}\heiti } \\[0.5\baselineskip]

    {\zihao{4}\kaishu \textsuperscript{1}，\textsuperscript{2}} \\[0.3\baselineskip]

    {\zihao{-5}\kaishu (1. ；2. )} \\[0.5\baselineskip]
\end{center}

% 中文摘要和关键词 - 左右各缩进2字符
\begin{quote}
    \zihao{-5}\songti
    \setlength{\leftskip}{2em}
    \setlength{\rightskip}{2em}

    \noindent{\heiti 摘\quad 要：}

    \vspace{0.5\baselineskip}
    \noindent{\heiti 关键词：}

    \noindent{\heiti 中图分类号：}\quad{\heiti 文献标识码：}A
\end{quote}

\vspace{1cm}

% 英文标题部分 - 通栏排版
\begin{center}
    {\fontsize{15}{18}\selectfont\bfseries } \\[0.3\baselineskip]

    {\fontsize{10.5}{12.6}\selectfont \textsuperscript{1}, \textsuperscript{2}} \\[0.3\baselineskip]

    {\fontsize{9}{10.8}\selectfont\itshape (1. ; 2. )} \\[0.5\baselineskip]
\end{center}

% 英文摘要和关键词
\begin{quote}
    \fontsize{9}{10.8}\selectfont
    \setlength{\leftskip}{2em}
    \setlength{\rightskip}{2em}

    \noindent{\bfseries Abstract：}
    
    \vspace{0.5\baselineskip}
    \noindent{\bfseries Keywords：}
\end{quote}

\vspace{1cm}

% 正文开始 - 双栏排版
\begin{multicols}{2}

    % 引言
    \section*{\zihao{-4}\heiti 引言}
    \zihao{5}\songti
    \setlength{\parindent}{2em}

    \section{\zihao{-4}\heiti 实验原理与方法}
    \subsection{\zihao{5}\heiti }
    \zihao{5}\songti

\subsection{\zihao{5}\heiti }

    \subsection{\zihao{5}\heiti }

    \section{\zihao{-4}\heiti 结果与讨论}
    \subsection{\zihao{5}\heiti }
    \zihao{5}\songti

      \subsection{\zihao{5}\heiti }

\section{\zihao{-4}\heiti 结\quad 论}
    \zihao{5}\songti

\end{multicols}

% 参考文献 - 通栏排版
%\section*{\zihao{-4}\heiti 参考文献}
%\zihao{-5}\songti

%\begin{enumerate}
%    \item 作者. 文献题名[J]. 刊名，出版年，卷(期): xxx-xxx（起止页码）.
%    \item 作者. 析出文献题名[C]//论文集主要责任者. 论文集名. 出版地：出版者，出版年.
%    \item 作者. 书名[M]. 版本(第一版不写). 出版地：出版者，出版年.
%    \item 作者. 文献题名[D]. 保存地点: 保存单位, 出版年.
%    \item 作者. 文献题名[EB/OL]. 电子文献的出处或可获得地址，发表或更新日期/引用日期.
%    \item 专利所有者. 专利题名[P]. 专利国别: 专利号, 出版日期.
%\end{enumerate}

\end{document}